
\section{无标记点圆盘模空间的定向}
\label{sec:8-1}

\subsection{全纯圆盘情形}
\label{subsec:8-1-1}

设 $(M,\omega)$ 为辛流形，$L$ 为拉格朗日子流形。
选取一个与 $\omega$ 相容的殆复结构 $J$。
于是 $L$ 是全实子流形，即 $J(TL)\cap TL$ 是 $TM|_L$ 的零截面。
本节考察带全实边界条件的伪全纯圆盘\liftIndex{I298}
\[
w\colon(D^2,\partial D^2)\longrightarrow(M,L)
\]
之模空间的定向问题。首先应当指出，伪全纯圆盘的模空间并不总是
可定向的；de Silva 在 \textbf{[Sil97]} 中也独立观察到了这一点。
另一方面，无边界伪全纯曲线的模空间众所周知带有典范定向。
为了讨论这里的定向问题，本章对第二施蒂费尔--惠特尼类
$w_2(TL)$ 作一个假设。本节的主要结果如下。

\begin{theorem}{8.1.1}\label{thm:8-1-1}
若 $L\subset(M,\omega)$ 是相对 spin 拉格朗日子流形，
则伪全纯圆盘的模空间可定向。此外，$L$ 上一个相对 spin 结构的选择，
对所有 $\beta\in\pi_2(M,L)$ 都典范地确定 $\mathcal M(L;\beta)$
上的一个定向。
\end{theorem}

下面给出相对 spin 结构的定义。设 $L\subset M$ 是定向拉格朗日
子流形，并设
\[
st\in H^2(M;\mathbb Z_2),\qquad st|_L=w_2(L).
\]
先固定 $M$ 的一个以 $L$ 为子复形的三角剖分。
在 $M$ 的 $3$-骨架 $M_{[3]}$ 上选取定向实向量丛 $V$，使得
$w_2(V)=st$。由于
\[
w_2\bigl(TL|_{L_{[2]}}\oplus V|_{L_{[2]}}\bigr)=0,
\]
$TL|_{L_{[2]}}\oplus V|_{L_{[2]}}$ 具有 spin 结构。

\begin{definition}{8.1.2}\label{def:8-1-2}
下列资料的选择称为拉格朗日子流形 $L\subset M$ 的一个
\emph{相对 spin 结构}\liftIndex{I365}：$L$ 的一个定向；$3$-骨架 $M_{[3]}$
上的一个定向实向量丛 $V$；一个上同调类
$st\in H^2(M;\mathbb Z_2)$；以及
$(TL\oplus V)|_{L_{[2]}}$ 上的一个 spin 结构 $\sigma$。
\end{definition}

若拉格朗日子流形 $L$ 带有相对 spin 结构，就称 $L$ 为
\emph{相对 spin 拉格朗日子流形}。

一对拉格朗日子流形 $(L^{(0)},L^{(1)})$ 称为\emph{相对 spin 的}，
如果存在一个类 $st$，同时满足
\[
st|_{L^{(i)}}=w_2(TL^{(i)})\qquad(i=0,1).
\]
该对子 $(L^{(0)},L^{(1)})$ 的一个相对 spin 结构，是指如下资料：
$L^{(i)}$ 的定向；$3$-骨架 $M_{[3]}$ 上满足 $st=w_2(V)$
的一个定向实向量丛 $V$；以及
\[
(TL^{(i)}\oplus V)|_{L^{(i)}_{[2]}}
\]
上的 spin 结构。这里 $L^{(i)}_{[2]}$ 是 $L^{(i)}$ 的 $2$-骨架。

\begin{remark}{8.1.3}\label{rem:8-1-3}
\begin{enumerate}[label=(\arabic*),leftmargin=2.8em]
\item 若 $L$ 是 spin 的，则可以如下地把 $L$ 的每个 spin 结构
对应到一个相对 spin 结构：取 $st=0$，令 $V$ 为平凡丛，
并在 $TL\oplus V=TL$ 上取由 $L$ 的 spin 结构诱导的 spin 结构。

\item 定义~\ref{def:8-1-2} 依赖于 $V$ 以及 $M$ 的三角剖分之选择。
稍后在命题~\ref{prop:8-1-6} 中，我们将在相对 spin 结构集合上
定义稳定共轭，使相对 spin 结构的稳定共轭类之集合不依赖于这些选择。
\end{enumerate}
\end{remark}

设 $D\colon\mathcal V\to\mathcal W$ 为弗雷德霍姆算子。
若行列式线
\[
\det D=\det(\operatorname{Coker}D)^*\otimes\det\operatorname{Ker}D
\]
带定向，就称 $D$ 的指标带定向。下面的事实将会反复使用。

选取 $\mathcal W$ 的有限维子空间 $E$，使
$\mathcal W=\operatorname{Im}D+E$。于是有典范同构
\[
\det D\cong\det E^*\otimes\det D^{-1}(E).
\]
对定向向量空间 $E$，我们如下规定 $E^*$ 的定向：若
$(e_1,\ldots,e_n)$ 是 $E$ 的定向基，
$(e_1^*,\ldots,e_n^*)$ 是 $E^*$ 的对偶基，则以
$(e_n^*,\ldots,e_1^*)$ 给出 $E^*$ 的定向。于是
\[
\det(S\oplus E)^*=\det E^*\otimes\det S^*.
\]
如下文所见，这一性质保证 $\det D$ 的定向不依赖于满足
$\mathcal W=\operatorname{Im}D+E$ 的有限维子空间 $E$ 之选择。
因此，若 $E$ 带有定向，例如带有复向量空间的定向，
则 $D^{-1}(E)$ 的定向确定 $\det D$ 的定向。

设 $E_1,E_2$ 是 $\mathcal W$ 的有限维子空间，并且
$\mathcal W=\operatorname{Im}D+E_i$（$i=1,2$）。
那么 $E_3=E_1+E_2$ 也满足这一性质。故只须证明：若
$E\subset E'$ 是满足 $\mathcal W=\operatorname{Im}D+E$ 的有限维
子空间，则定向线
\[
\det E^*\otimes\det D^{-1}(E),\qquad
\det E'^*\otimes\det D^{-1}(E')
\]
自然同构。为此，选取
$S\subset E'\cap\operatorname{Im}D$，使得作为 $\mathcal W$ 的
子空间有 $E'=S\oplus E$；再选取 $\mathcal V$ 的子空间
$\widetilde S$，使限制
\[
D|_{\widetilde S}\colon\widetilde S\longrightarrow S
\]
是有限维向量空间之间的同构。于是
\[
D^{-1}(S\oplus E)=\widetilde S\oplus D^{-1}(E)
\]
是 $\mathcal V$ 中有限维子空间的等式。因此
\[
\begin{split}
\det(S\oplus E)^*\otimes\det D^{-1}(S\oplus E)
={}&\det E^*\otimes\det S^*\\
&\otimes\det\widetilde S\otimes\det D^{-1}(E).
\end{split}
\]
由于 $D|_{\widetilde S}\colon\widetilde S\to S$ 是同构，
它诱导从 $\det S^*\otimes\det\widetilde S$ 到定向线
$\mathbb R$ 的典范同构。总之，
$\det E'^*\otimes\det D^{-1}(E')$ 与
$\det E^*\otimes\det D^{-1}(E)$ 作为定向线同构。%
\translatornote{源书在本句两处 $D^{-1}(E')$、$D^{-1}(E)$ 前漏写
$\det$；依紧前的完整公式及两侧必须是线这一型别补出。}

如上所述，$E$ 的定向诱导 $E^*$ 的定向。因此有时我们将考察
$E$ 的定向，而不直接考察 $E^*$ 的定向。

首先研究线性化问题中的定向。

\begin{proposition}{8.1.4}\label{prop:8-1-4}
考虑 $(D^2,\partial D^2)$ 上的复丛对 $(E,\lambda)$。
假设 $\lambda$ 平凡，则 $\lambda$ 的每个平凡化都典范地诱导
$\operatorname{Index}\bar\partial_{(E,\lambda)}$ 的一个定向。
\end{proposition}

\begin{proof}
固定 $\partial D^2$ 的一个环形邻域
$A(\epsilon)\subset D^2$，其中%
\translatornote{源书把本句写成“each $\partial D^2$”，并把环境空间
误写为未定义的 $D_i$；这里依紧接的集合定义恢复为唯一的圆盘 $D^2$。}
\[
A(\epsilon)=\{z\in D^2\mid 1-\epsilon\leq |z|\leq1\},
\qquad 0<\epsilon<1.
\]
记圆周
\[
C_\epsilon=S^1(1-\epsilon)\subset A(\epsilon).
\]
沿 $C_\epsilon$ 把 $D^2$ 收缩，并把 $C_\epsilon$ 识别为一点，
所得商空间 $D^2/\!\sim$ 具有节点曲线的结构。记商映射为
\[
\operatorname{cont}\colon D^2\longrightarrow
\Sigma:=D^2/\!\sim\ \cong\Sigma_{\mathrm{ct}}\cup D,
\]
其中 $\Sigma_{\mathrm{ct}}\cong\mathbb{CP}^1$ 且 $D\cong D^2$。
下文把 $\Sigma_{\mathrm{ct}}$ 简记为 $\mathbb{CP}^1$。
以 $p$ 表示 $\mathbb{CP}^1$ 上与 $D$ 中原点 $0$ 相接的节点，
并置 $S=\partial D$。

\begin{figure}[!htbp]
 \centering
 \setlength{\unitlength}{0.72mm}
 \begin{picture}(118,58)
  \linethickness{0.55pt}
  \put(27,29){\circle{36}}
  \linethickness{0.30pt}
  {\color{black!55}\put(27,29){\circle{25}}}
  \put(17,43){\makebox(0,0)[l]{$\partial D^2$}}
  \put(28,18){\makebox(0,0){$C_\epsilon$}}
  \linethickness{0.55pt}
  \put(52,29){\vector(1,0){15}}
  \put(52,31){\line(1,0){12}}
  \put(88,22){\oval(39,25)}
  \put(88,22){\circle*{1.7}}
  \qbezier(88,22)(76,44)(87,50)
  \qbezier(87,50)(101,47)(88,22)
  \put(90,19){\makebox(0,0)[l]{$C_\epsilon$}}
 \end{picture}
 \par\smallskip
 \renewcommand{\thefigure}{8.1.1}
 \refstepcounter{figure}
 \textbf{图~\thefigure}\label{fig:8-1-1}
\end{figure}

利用给定的同构
\[
E|_{\partial D^2}\cong\lambda|_{\partial D^2}\otimes\mathbb C,
\]
$\lambda\to\partial D^2$ 的每个平凡化都诱导复平凡化
\[
\Phi\colon E|_{A(\epsilon)}\longrightarrow A(\epsilon)\times\mathbb C^n.
\]
因此 $D^2$ 上的丛 $E$ 下降为节点曲线
$\mathbb{CP}^1\cup D$ 上的丛对 $(E',\lambda')$。

先研究 $(E',\lambda')$ 的定向问题。这个丛对的一个截面是截面对
$(\xi_{\mathbb{CP}^1},\xi_D)$，其中
$\xi_{\mathbb{CP}^1}$ 是 $E'|_{\mathbb{CP}^1}$ 的截面，
$\xi_D$ 是 $E'\to D$ 的截面，并满足
\[
\xi_D(z)\in\lambda'_z\quad(z\in S),
\qquad
\xi_{\mathbb{CP}^1}(p)=\xi_D(0).
\]
换言之，$(E',\lambda')$ 的 $W^{1,p}$-截面空间为
\[
W^{1,p}(E',\lambda'):=\operatorname{diff}^{-1}(0),
\]
其中
\[
\begin{split}
\operatorname{diff}\colon{}
W^{1,p}(E'|_{\mathbb{CP}^1})\oplus
W^{1,p}(E'|_D,\lambda')&\longrightarrow\mathbb C^n,\\
(\xi_{\mathbb{CP}^1},\xi_D)&\longmapsto
\xi_{\mathbb{CP}^1}(p)-\xi_D(0).
\end{split}
\]
这给出正合列
\[
\begin{split}
0\longrightarrow W^{1,p}(E',\lambda')
\longrightarrow{}&W^{1,p}(E'|_{\mathbb{CP}^1})
 \oplus W^{1,p}(E'|_D,\lambda')\\
\longrightarrow{}&\mathbb C^n\longrightarrow0.
\end{split}
\]
丛对 $(E',\lambda')$ 在 $\Sigma=D^2/\!\sim$ 上的相应指标，
是算子
\[
\begin{split}
\bar\partial_{(E',\lambda')}\colon
W^{1,p}(E',\lambda')\longrightarrow{}
&L^p\bigl(\Lambda^{(0,1)}\otimes E'|_{\mathbb{CP}^1}\bigr)\\
&\oplus L^p\bigl(\Lambda^{(0,1)}\otimes E'|_D\bigr)
\end{split}
\]
的指标。

下面证明：$\lambda\to\partial\Sigma=\partial D$ 的每个平凡化，
都典范地确定 $\operatorname{Index}\bar\partial_{(E',\lambda')}$
上的一个定向，从而也确定
$\det\bar\partial_{(E',\lambda')}$ 上的定向。
由 $\Sigma$ 与其上丛对 $(E',\lambda')$ 的构造，
$\lambda\to\partial D^2$ 的平凡化典范地诱导
$\lambda'\to\partial\Sigma$ 的平凡化。

首先注意，
\[
\begin{split}
\bar\partial_{(E'|_D,\lambda'|_S)}\colon
W^{1,p}(E'|_D,\lambda'|_S)
\longrightarrow L^p(\Lambda^{(0,1)}\otimes E'|_D)
\end{split}
\]
是满射，并且
\[
\ker\bar\partial_{(E'|_D,\lambda'|_S)}
\cong\ker\bar\partial_{(D\times\mathbb C^n,S\times\mathbb R^n)}
\cong\mathbb R^n.
\]
事实上，$\ker\bar\partial_{(E'|_D,\lambda'|_S)}$ 的每个元素，
都是柯西--黎曼方程
\[
\left\{
\begin{aligned}
\bar\partial\xi&=0,\\
\xi(z)&\in\lambda'_z\equiv\mathbb R^n
\end{aligned}
\right.
\]
的解 $\xi\colon D\cong D^2\to\mathbb C^n$，故必为实常向量。
这证明了典范同构
\[
\ker\bar\partial_{(E'|_D,\lambda'|_S)}\cong\mathbb R^n.
\]

于是可以选取一个有限维复子空间
\[
\mathcal E\subset
L^p(\Lambda^{(0,1)}\otimes E'|_{\mathbb{CP}^1}),
\qquad \dim_{\mathbb C}\mathcal E=k,
\]
使得
\[
\begin{split}
(\mathcal E\oplus0)+\operatorname{Range}
 \bar\partial_{(E',\lambda')}
={}&L^p(\Lambda^{(0,1)}\otimes E'|_{\mathbb{CP}^1})\\
&\oplus L^p(\Lambda^{(0,1)}\otimes E'|_D),
\end{split}
\]
且评价映射
\[
\operatorname{ev}_p\colon
\bigl(\bar\partial_{E'|_{\mathbb{CP}^1}}\bigr)^{-1}(\mathcal E)
\longrightarrow E'_p\cong\mathbb C^n
\]
为满射。固定一个同构
\[
\Psi\colon\mathbb C^k\longrightarrow\mathcal E\oplus0,
\]
并考虑算子
\[
\widehat{\bar\partial}_{(E',\lambda');\Psi}(h,\xi)
:=\Psi(h)+\bar\partial_{(E',\lambda')}(\xi)
\]
以及算子
\[
\begin{split}
\bar\partial_{(E',\lambda');\Psi}\colon
\mathbb C^k\oplus W^{1,p}(E',\lambda')
\longrightarrow{}&\mathbb C^k\oplus
L^p(\Lambda^{(0,1)}\otimes E'|_{\mathbb{CP}^1})\\
&\oplus L^p(\Lambda^{(0,1)}\otimes E'|_D),
\end{split}
\]
后者定义为
\[
\bar\partial_{(E',\lambda');\Psi}(h,\xi)
=\bigl(0,\widehat{\bar\partial}_{(E',\lambda');\Psi}(h,\xi)\bigr).
\]
由 $\mathcal E$ 的选取，
$\widehat{\bar\partial}_{(E',\lambda');\Psi}$ 为满射。

我们有正合列
\[
\begin{split}
0\longrightarrow\ker\bar\partial_{(E',\lambda')}
\longrightarrow\ker\bar\partial_{(E',\lambda');\Psi}
\longrightarrow\mathbb C^k
\longrightarrow\operatorname{Coker}\bar\partial_{(E',\lambda')}
\longrightarrow0,
\end{split}
\]
其中依次位于中间的三个映射为
\[
\begin{split}
d_1(\xi)&=(0,\xi),\\
d_2(h,\xi)&=h,\\
d_3(h)&=\Psi(h)\pmod{\operatorname{Range}\bar\partial_{(E',\lambda')}}.
\end{split}
\]
还要注意，$\bar\partial_{(E',\lambda');\Psi}$ 通过弗雷德霍姆算子
同伦于 $\operatorname{id}_{\mathbb C^k}\oplus
\bar\partial_{(E',\lambda')}$。因此得到
\begin{equation}
\operatorname{Index}\bar\partial_{(E',\lambda')}
\cong\operatorname{Index}\bar\partial_{(E',\lambda');\Psi}
\cong\ker\widehat{\bar\partial}_{(E',\lambda');\Psi}-\mathbb C^k.
\tag{*}\label{eq:8-1-star}
\end{equation}

注意 $\ker\widehat{\bar\partial}_{(E',\lambda');\Psi}$ 是满足
\[
0=\Psi(h)+\bar\partial_{(E',\lambda')}(\xi)
\]
的对子 $(h,\xi)$ 之集合，即
\[
\ker\widehat{\bar\partial}_{(E',\lambda');\Psi}
\cong\bigl(\bar\partial_{(E',\lambda')}\bigr)^{-1}
(0\oplus\mathcal E).
\]
由于 $\mathcal E$ 的选取使评价映射
\[
\operatorname{ev}_p\colon
\bigl(\bar\partial_{E'|_{\mathbb{CP}^1}}\bigr)^{-1}(\mathcal E)
\longrightarrow E'_p\cong\mathbb C^n
\]
为满射，故
$\bigl(\bar\partial_{(E',\lambda')}\bigr)^{-1}
(0\oplus\mathcal E)$ 正是下列满同态的核：
\[
\begin{split}
\operatorname{diff}\colon{}
&\bigl(\bar\partial_{E'|_{\mathbb{CP}^1}}\bigr)^{-1}(\mathcal E)
\times\ker\bar\partial_{(E'|_D,\lambda'|_S)}
\longrightarrow\mathbb C^n,\\
&(\xi_{\mathbb{CP}^1},\xi_D)
\longmapsto\xi_{\mathbb{CP}^1}(p)-\xi_D(0).
\end{split}
\]
回忆，
$\bigl(\bar\partial_{E'|_{\mathbb{CP}^1}}\bigr)^{-1}(\mathcal E)$
和 $\mathbb C^n$ 都是复向量空间，因而典范地带定向。

这说明，给定 $\lambda$ 的一个平凡化以后，它典范地确定
$\ker\bar\partial_{(E'|_D,\lambda'|_S)}$ 的定向，
从而确定
$\bigl(\bar\partial_{(E',\lambda')}\bigr)^{-1}
(0\oplus\mathcal E)$ 的定向。再由同构~\eqref{eq:8-1-star}，
它确定 $\operatorname{Index}\bar\partial_{(E',\lambda')}$ 的定向。

下面通过粘合过程把这个定向传到 $(E,\lambda)$ 上。
可以磨光节点，得到由 $r\in\mathbb R_+$ 参数化的一族黎曼曲面，
使丛对 $(E',\lambda')$ 本身形变为
$(\Sigma_r,\partial\Sigma_r)$ 上的丛对 $(E_r,\lambda_r)$，
并有典范识别 $\lambda'\cong\lambda_r$。

更具体地，在 $0\in D$ 与 $p\in\mathbb{CP}^1$ 附近分别选取
小坐标邻域 $D_\delta(0)$ 与 $D_\delta(p)$；其中
$D_\delta(0)$ 是以 $0$ 为中心、半径为 $\delta$ 的圆盘。
对正实数 $r>0$，在 $0$ 与 $p$ 附近粘合 $D^2$ 和
$\mathbb{CP}^1$：当
\[
z\in D_\delta(0),\qquad w\in D_\delta(p),\qquad zw=1/r
\]
时识别 $z$ 与 $w$。%
\translatornote{源书先把节点两支记为 $0,p$，下一句却改写成未定义于此的
$O,S$；其中 $S$ 此前已表示整条边界 $\partial D$。这里依节点定义统一恢复
为 $0,p$。}
记所得有边黎曼曲面为 $\Sigma_r$；它双全纯于单位圆盘。
我们可以识别
\[
\Sigma_1=D^2,
\qquad
\Sigma_\infty=\mathbb{CP}^1\cup D.
\]
从 $E'$ 也得到 $\Sigma_r$ 上的向量丛，记作 $E_r$；
$\partial D$ 上的全实子丛 $\lambda'$ 诱导
$\partial\Sigma_r$ 上的全实子丛 $\lambda_r$。

由于 $\mathcal E$ 中元素的支集都避开节点，$\mathcal E$ 典范地嵌入
\[
L^p\bigl(\Lambda^{(0,1)}(\Sigma_r)\otimes E_r\bigr),
\]
所以当 $r$ 充分大时，可以典范地把 $\Psi$ 看作映射
\[
\Psi\colon\mathbb C^k\longrightarrow
L^p\bigl(\Lambda^{(0,1)}(\Sigma_r)\otimes E_r\bigr).
\]
利用 $\widehat{\bar\partial}_{(E',\lambda');\Psi}$ 的满射性以及
粘合理论（比较 \textbf{[SeSi88]}），有
\[
\operatorname{Range}
\widehat{\bar\partial}_{(E_r,\lambda_r);\Psi}+\mathcal E
=L^p\bigl(\Lambda^{(0,1)}(\Sigma_r)\otimes E_r\bigr)
\]
以及典范同构
\[
\operatorname{Index}\bar\partial_{(E',\lambda');\Psi}
\cong
\operatorname{Index}\bar\partial_{(E_r,\lambda_r);\Psi}.
\]
所以，当 $r\to\infty$ 时
$\operatorname{Index}\bar\partial_{(E',\lambda');\Psi}$ 上的定向，
典范地给出充分大 $r>0$ 时
$\operatorname{Index}\bar\partial_{(E_r,\lambda_r);\Psi}$ 上的定向；
再由同构~\eqref{eq:8-1-star} 的第一部分，得到
$\operatorname{Index}\bar\partial_{(E_r,\lambda_r)}$ 上的定向。
最后如上所述，把 $(E_r,\lambda_r)$ 连续形变到原先给定的丛对
$(E_1,\lambda_1)=(E,\lambda)$，便在后者上诱导出定向。
证明完毕。
\end{proof}

\begin{proof}[定理~\ref{thm:8-1-1} 的证明]
把命题~\ref{prop:8-1-4} 应用于
\[
E=w^*TM,
\qquad
F=\ell^*TL,
\qquad
\ell=w|_{\partial D^2}.
\]
由于 $L$ 已定向，$\ell^*TL$ 是平凡丛；但它的平凡化并不唯一，
而一个相对 spin 结构的选择会给出一个在同伦意义下唯一的
（稳定）平凡化。

先概述证明。利用命题~\ref{prop:8-1-4}，对每个全纯圆盘 $w$
赋予 $\det D_w\bar\partial$ 的一个定向，然后证明该定向只依赖于
相对 spin 结构。具体步骤如下。首先，在
$(D^2,\partial D^2)\to(M,L)$ 的每个映射同伦类中选取一个全纯圆盘
$w_0$，并定义 $\det D_{w_0}\bar\partial$ 的定向；这个定向依赖于
相对 spin 结构之选择。利用下面的命题 8.1.7 可知，
若改变 $\ell^*TL$ 的稳定平凡化之同伦类，则
命题~\ref{prop:8-1-4} 证明中指标丛的定向也会改变。

具体而言，设 $g$ 是 $\ell^*TL$ 的一个规范变换，它交换稳定平凡化
的两个同伦类。由于 $SO(n)$ 中的任意环路在 $U(n)$ 中可缩，
可以把 $g\otimes\mathbb C$ 延拓为 $w^*TM$ 的一个规范变换
$\widetilde g$。考虑
\[
(w^*TM,\ell^*TL)\times[0,1].
\]
利用 $\widetilde g$ 识别
$(w^*TM,\ell^*TL)\times\{0\}$ 与
$(w^*TM,\ell^*TL)\times\{1\}$，便得到
$(D^2,\partial D^2)\times S^1$ 上的一个丛对。
这个丛对稳定等价于下一小节的 $(E,F)$，所以与
$\widetilde g$ 相应的丛对之族指标不可定向。
这说明，改变 $\ell^*TL$ 的稳定平凡化会改变指标丛的定向。

若 $\{w_t\}_{0\leq t\leq1}$ 是从 $w_0$ 到 $w=w_1$ 的同伦，
则行列式线丛 $\det D_{w_t}\bar\partial$ 是平凡的，
因为底空间 $[0,1]$ 可缩。因此得到
$\det D_w\bar\partial$ 的一个定向，但它或许依赖于同伦
$\{w_t\}_{0\leq t\leq1}$。余下任务是证明所诱导的定向不依赖于
同伦之选择。这里需要 $L$ 相对 spin，也就是说
$L\subset M$ 上存在相对 spin 结构。（下面第 2 步证明 $S^1$ 上某个
行列式线丛可定向；这一步只要求相对 spin 结构存在，
而不依赖于某个特定选择。）

\medskip
\noindent\textbf{第 1 步：为 $\det D_w\bar\partial$ 赋予定向。}
记
\[
\begin{split}
D_w\bar\partial\colon
W^{1,p}(D^2,\partial D^2;w^*TM,\ell^*TL)
\longrightarrow{}&W^{0,p}(D^2;w^*TM\otimes\Lambda^{0,1}(D^2))\\
={}&L^p(D^2;w^*TM\otimes\Lambda^{0,1}(D^2))
\end{split}
\]
为伪全纯曲线方程的线性化算子；它是一阶椭圆微分算子，
主符号与多尔博算子相同。这里 $\ell$ 是 $w$ 在边界
$\partial D^2$ 上的限制，并取 $p>2$。
要证明定理~\ref{thm:8-1-1}，只须证明线性化算子
$D_w\bar\partial$ 的指标带定向。

注意，对任意零阶算子 $A$，$D_w\bar\partial+A$ 总是
弗雷德霍姆算子；所有这类算子组成的空间
$\mathcal F_{D_w\bar\partial}$ 可缩。
族 $D_w\bar\partial+A\in\mathcal F_{D_w\bar\partial}$ 的指标，
是这个可缩空间上的虚向量丛；特别地，该族的行列式线丛平凡。
零阶项不影响指标问题，所以可以用多尔博算子
\[
\begin{split}
\bar\partial_{(w^*TM,\ell^*TL)}\colon
W^{1,p}(D^2,\partial D^2;w^*TM,\ell^*TL)
\longrightarrow L^p(D^2;w^*TM\otimes\Lambda^{0,1}(D^2))
\end{split}
\]
代替 $D_w\bar\partial$。

由单纯逼近定理，任意映射
$w\colon(D^2,\partial D^2)\to(M,L)$ 都可以形变为某个映射
\[
w_0\colon(D^2,\partial D^2)\longrightarrow(M_{[2]},L_{[1]}).
\]
（在映射空间
$(D^2,\partial D^2)\to(M,L)$ 的每个连通分支中选取一个这样的
$w_0$。）相对 spin 结构确定
$\ell_0^*(TL\oplus V)$ 的平凡化之稳定同伦类。
由命题~\ref{prop:8-1-4}，可以指定
\[
\det\bigl(\operatorname{Coker}
 \bar\partial_{(w_0^*TM,\ell_0^*TL)}\bigr)^*
\otimes
\det\ker\bar\partial_{(w_0^*TM,\ell_0^*TL)}
\]
的一个定向。

对包含 $w_0$ 的映射空间道路分支中任意给定的 $w$，
考虑一条从 $w_0$ 出发、终止于 $w$，并由 $t\in[0,1]$ 参数化的
映射道路
\[
w_t\colon(D^2,\partial D^2)\longrightarrow(M,L).
\]
（一旦在 $TM$ 上选定复线性联络，线性算子
$D_w\bar\partial$ 对任意光滑映射
$(D^2,\partial D^2)\to(M,L)$ 都有意义。）
由命题~\ref{prop:8-1-4} 及 $[0,1]$ 的可缩性，
在 $S^1\times[0,1]$ 上取 $\ell_t^*TL$ 的一个平凡化，其中
$\ell_t=w_t|_{\partial D^2}$，便可为
\[
\det\bigl(\operatorname{Coker}
 \bar\partial_{(w_t^*TM,\ell_t^*TL)}\bigr)^*
\otimes
\det\ker\bar\partial_{(w_t^*TM,\ell_t^*TL)}
\]
给出一个随 $t\in[0,1]$ 连续变化的定向族。
这就为 $w$ 定义了一个定向。

\medskip
\noindent\textbf{第 2 步：与道路 $\{w_t\}$ 的选择无关。}
我们证明上述定向的赋值不依赖于道路
$\{w_t\}_{t\in[0,1]}$ 之选择；这里要用到相对 spin 结构的存在性。
设 $st,V$ 如定义~\ref{def:8-1-2}。
相对 spin 结构在 $M_{[2]}\cap L$ 上确定
$TL\oplus V$ 的一个 spin 结构 $\sigma$。

我们要证明，从 $w_0$ 到 $w$ 的任意两条道路
$w_t^{(1)},w_t^{(2)}$ 都诱导相同的指标定向。
把 $w_t^{(1)}$ 与反向的 $w_t^{(2)}$ 连接起来，得到一个由 $S^1$
参数化的映射族 $(D^2,\partial D^2)\to(M,L)$。
更一般地，设
\[
w_t\colon(D^2,\partial D^2)\longrightarrow(M,L)
\qquad(t\in S^1)
\]
是这样的一族映射，并定义
\[
\Phi\colon D^2\times S^1\longrightarrow M,
\qquad
\Phi(z,t)=w_t(z).
\]
（特别地，下面的论证适用于由 $S^1$ 参数化的伪全纯圆盘族。）

由单纯逼近定理，可以把 $\Phi$ 同伦为映射 $\Phi'$，使得
\[
\Phi'(D^2\times S^1)\subset M_{[3]},
\qquad
\Phi'(\partial D^2\times S^1)
\subset M_{[2]}\cap L=L_{[2]}.
\]
与 $\Phi'$ 对应的、由 $S^1$ 参数化的多尔博型算子族，
在弗雷德霍姆算子空间中同伦于与 $\Phi$ 对应的算子族。
所以只要讨论指标，尤其是定向问题，只须考虑 $\Phi'$。
以下置 $\Phi'=\Phi$。

把 $V$ 由 $\Phi$ 拉回到 $D^2\times S^1$ 上，记作 $\Phi^*V$。
在 $(TL\oplus V)|_{L_{[2]}}$ 上选定的 spin 结构给出它在
$L_{[2]}$ 上的稳定平凡化，进而诱导
\[
\Phi^*(TL\oplus V)|_{\partial D^2\times S^1}
\]
的稳定平凡化。另一方面，$\Phi^*V$ 定义在
$D^2\times S^1\simeq S^1$ 上且已经定向。
由于底空间是 $D^2\times S^1$，$\Phi^*V$ 的定向保证它在
$D^2\times S^1$ 上存在平凡化，因而其限制在
$\partial D^2\times S^1$ 上也有平凡化。
还要指出，这个平凡化在 $\partial D^2\times\{t\}$ 上的限制，
不依赖于 $(TL\oplus V)|_{L_{[2]}}$ 上 spin 结构之选择。%
\translatornote{源书此处写成 $(TL\oplus V)|_{M_{[3]}\cap L}$，
但定义 8.1.2 及本步开头均只在 $M_{[2]}\cap L=L_{[2]}$ 上给出
该 spin 结构；据此恢复定义域。}

合并上述两个平凡化，便得到
\[
\bigl(\Phi|_{\partial D^2\times\{t\}}\bigr)^*TL
\]
的一族稳定平凡化，它随 $t\in S^1$ 连续变化。
特别地，$\ell^*(TL\oplus V)$ 的平凡化之同伦类不依赖于
道路 $w_t^{(i)}$（$i=1,2$）之选择。综上，对任意 $w$，
$(TL\oplus V)|_{L_{[2]}}$ 上的 spin 结构诱导
$\ell^*(TL\oplus V)$ 的一个唯一平凡化同伦类，
其中 $\ell=w|_{\partial D^2}$；而 $M_{[3]}$ 上定向丛 $V$ 的存在，
诱导 $\ell^*V$ 的一个唯一平凡化同伦类。
也就是说，对每个 $\ell=w|_{\partial D^2}$，$\ell^*TL$
都得到稳定平凡化；而且可以让这个稳定平凡化随 $\ell$ 连续变化。

置 $V_{\mathbb C}=V\otimes\mathbb C$。于是
\[
\begin{aligned}
&\bar\partial_{(w^*(TM\oplus V_{\mathbb C}),
 \ell^*(TL\oplus V))}\colon\\[-0.2em]
&\quad
W^{1,p}\bigl(w^*(TM\oplus V_{\mathbb C}),
 \ell^*(TL\oplus V)\bigr)
\longrightarrow
L^p\bigl(w^*(TM\oplus V_{\mathbb C})\bigr)
\end{aligned}
\]
是下列两个算子的直和：
\[
\bar\partial_{(w^*TM,\ell^*TL)}\colon
W^{1,p}(w^*TM,\ell^*TL)\longrightarrow L^p(w^*TM)
\]
以及
\[
\bar\partial_{(w^*V_{\mathbb C},\ell^*V)}\colon
W^{1,p}(w^*V_{\mathbb C},\ell^*V)
\longrightarrow L^p(w^*V_{\mathbb C}).
\]
由命题~\ref{prop:8-1-4}，$\ell^*(TL\oplus V)$ 的定向典范地诱导
$\bar\partial_{(w^*(TM\oplus V_{\mathbb C}),
\ell^*(TL\oplus V))}$ 的指标定向。
$\ell^*V$ 的平凡化延拓到 $w^*V$，因而典范地诱导
$\bar\partial_{(w^*V_{\mathbb C},\ell^*V)}$ 的定向。
所以 $\bar\partial_{(w^*TM,\ell^*TL)}$ 的指标也典范地带定向，
从而 $\mathcal M(L;\beta)$ 带定向。

由上述论证，$L$ 的定向、类 $st\in H^2(M;\mathbb Z_2)$、
$M_{[3]}$ 上的向量丛 $V$，以及
$(TL\oplus V)|_{M_{[2]}\cap L}$ 上的 spin 结构之选择，
对所有 $\beta$ 都典范地确定 $\mathcal M(L;\beta)$ 的定向。
定理~\ref{thm:8-1-1} 得证。
\end{proof}

在后文的论证中，定理~\ref{thm:8-1-1} 证明第 1 步所给的定向
与全纯圆盘的粘合过程相容，这一点十分重要。

下面澄清相对 spin 结构的概念怎样依赖于 $V$ 和三角剖分之选择。
首先回忆一个基本事实。设 $E$ 是 $X$ 上满足 $w_2(E)=0$ 的定向
向量丛。一个 spin 结构等价于相应主 $SO$-丛 $P_E$ 的逐纤维二重覆盖；
它由同态
\[
\pi_1(P_E)\longrightarrow\mathbb Z_2
\]
刻画，并且该同态在纤维上的限制
$\pi_1(SO)\to\mathbb Z_2$ 非平凡。因此，$X$ 上 $E$ 的 spin 结构
之集合是 $H^1(X;\mathbb Z_2)$ 的一个挠子。

$E$ 上的 spin 结构会在其稳定化 $E\oplus\mathbb R^k$ 上诱导
spin 结构，其中 $\mathbb R^k$ 是带给定平凡化的向量丛；
所诱导的结构来自 $E$ 上的 spin 结构以及 $\mathbb R^k$ 的平凡化。
若只把 $\mathbb R^k$ 看作平凡丛而不指定具体平凡化，
就在稳定化后的 spin 结构之间得到一个粗等价关系；这些等价类由
$H^1(X;\mathbb Z_2)$ 对 $\mathbb R^k$ 的规范变换取商后参数化。
（比较上面对 spin 结构的描述。）现在要在相对 spin 结构集合上
引入类似的粗等价关系。

下一个定义中仍固定 $M$ 的一个三角剖分 $\mathcal T$，使 $L$ 是其
子复形；也就是说，$\mathcal T$ 与 $L$ 相容。

\begin{definition}{8.1.5}\label{def:8-1-5}
设 $(st_i,V_i,\sigma_i)$（$i=1,2$）是 $L$ 的相对 spin 结构。
如果存在整数 $k_i$ 以及保定向丛同构
\[
\tau\colon V_1\oplus\mathbb R^{k_1}
\longrightarrow V_2\oplus\mathbb R^{k_2},
\]
使得在
\[
\begin{split}
1\oplus\tau|_{L_{[2]}}\colon
(TL\oplus V_1)|_{L_{[2]}}\oplus\mathbb R^{k_1}
\longrightarrow
(TL\oplus V_2)|_{L_{[2]}}\oplus\mathbb R^{k_2}
\end{split}
\]
之下，spin 结构 $\sigma_1\oplus1$ 诱导 spin 结构
$\sigma_2\oplus1$，就称这两个相对 spin 结构\emph{稳定共轭}。
\end{definition}

注意，若 $(st_1,V_1,\sigma_1)$ 与 $(st_2,V_2,\sigma_2)$ 稳定共轭，
则它们在 $\mathcal M(\beta)$ 上确定相同的定向。
（这从定理~\ref{thm:8-1-1} 证明中给出的定向定义立即可见。）

对给定的三角剖分 $\mathcal T$，以
\[
\operatorname{Spin}(M,L;\mathcal T)
\]
表示相对于 $\mathcal T$ 的、$L$ 上所有相对 spin 结构之稳定共轭类
所成的集合。下面命题~\ref{prop:8-1-6}(2) 说明：对任意两个与 $L$
相容的三角剖分 $\mathcal T,\mathcal T'$，集合
$\operatorname{Spin}(M,L;\mathcal T)$ 与
$\operatorname{Spin}(M,L;\mathcal T')$ 典范地识别。
因此可以定义 $L$ 上相对 spin 结构的稳定共轭类之集合，记作
$\operatorname{Spin}(M,L)$。

\begin{proposition}{8.1.6}\label{prop:8-1-6}
\begin{enumerate}[label=(\arabic*),leftmargin=2.8em]
\item $\operatorname{Spin}(M,L;\mathcal T)$ 是
$H^2(M,L;\mathbb Z_2)$ 的一个挠子。

\item 对 $M$ 上任意两个与 $L$ 相容的三角剖分
$\mathcal T,\mathcal T'$，存在与上述作用相容的典范同构
\[
\operatorname{Spin}(M,L;\mathcal T)
\cong\operatorname{Spin}(M,L;\mathcal T').
\]
\end{enumerate}
\end{proposition}

特别地，若给定 $L$ 的一个 spin 结构，则有典范同构
\[
\operatorname{Spin}(M,L;\mathcal T)
\cong H^2(M,L;\mathbb Z_2).
\]

\begin{proof}
取
\[
[(st,V,\sigma)]\in\operatorname{Spin}(M,L;\mathcal T),
\qquad
\mathfrak x\in H^2(M,L;\mathbb Z_2).
\]
利用同构
\[
H^2(M,L;\mathbb Z_2)
\cong H^2(M_{[3]},L_{[2]};\mathbb Z_2),
\]
把 $\mathfrak x$ 看作右端中的类。由于
$M_{[3]}/L_{[2]}$ 是三维胞腔复形，映射
\[
\begin{split}
(w_1,w_2)\colon
\widetilde{KO}(M_{[3]}/L_{[2]})
\longrightarrow{}&H^1(M_{[3]},L_{[2]};\mathbb Z_2)\\
&\oplus H^2(M_{[3]},L_{[2]};\mathbb Z_2)
\end{split}
\]
为同构。因此 $\mathfrak x\in H^2(M,L;\mathbb Z_2)$ 唯一确定
$M_{[3]}/L_{[2]}$ 上一个可定向向量丛 $E_{\mathfrak x}$ 的稳定类，
并且
\[
w_2(E_{\mathfrak x})=\mathfrak x.
\]
选定 $E_{\mathfrak x}$ 的一个定向。
沿 $M_{[3]}\to M_{[3]}/L_{[2]}$ 拉回 $E_{\mathfrak x}$，
得到向量丛 $V_{\mathfrak x}$。注意，
$V_{\mathfrak x}|_{L_{[2]}}$ 带有由塌缩点处标架拉回的平凡化，
从而带有 spin 结构 $\sigma_{\mathfrak x}$。现在规定
\[
\mathfrak x\cdot[(st,V,\sigma)]
=\bigl[(st+\overline{\mathfrak x},
 V\oplus V_{\mathfrak x},
 \sigma\oplus\sigma_{\mathfrak x})\bigr],
\]
其中 $\overline{\mathfrak x}$ 是 $\mathfrak x$ 在映射
\[
H^2(M,L;\mathbb Z_2)\longrightarrow H^2(M;\mathbb Z_2)
\]
下的像。右端不依赖于为 $E_{\mathfrak x}$ 选取的定向，
因为 $E_{\mathfrak x}\oplus\mathbb R$ 上存在反转定向的对合，
即在第二个因子 $\mathbb R$ 上乘以 $-1$。

反过来，对两个相对 spin 结构
$[(st,V,\sigma)]$ 与 $[(st',V',\sigma')]$，如下定义它们的差。
在 $L_{[2]}$ 上选取定向向量丛 $E$，使 $TL\oplus E$ 平凡，
并固定同构
\[
I\colon TL\oplus E\xrightarrow{\cong}L_{[2]}\times\mathbb R^N.
\]
对 $V'$，在 $M_{[3]}$ 上选取定向向量丛 $W'$，使
$V'\oplus W'$ 平凡，并固定同构
\[
I_{V',W'}\colon V'\oplus W'\xrightarrow{\cong}
M_{[3]}\times\mathbb R^{N'}.
\]
spin 结构 $\sigma'$ 与同构 $I,I_{V',W'}$ 在
$E\oplus W'$ 上诱导一个 spin 结构，记作 $\check\sigma'$。
于是向量丛 $V\oplus W'$ 在 $L_{[2]}$ 上带有如下稳定平凡化：
\[
\begin{split}
(V\oplus W')|_{L_{[2]}}\oplus\mathbb R^N
&\cong (V\oplus W')|_{L_{[2]}}\oplus(TL\oplus E)\\
&\cong (TL\oplus V|_{L_{[2]}})
 \oplus(E\oplus W'|_{L_{[2]}}).
\end{split}
\]
这一平凡化与 spin 结构 $\sigma\oplus\check\sigma'$ 相容。
记它为 $\mathfrak I_{\sigma,\sigma'}$，并由之得到元素
\[
[V\oplus W',\mathfrak I_{\sigma,\sigma'}]
\in\widetilde{KO}(M_{[3]}/L_{[2]}).
\]

固定一个相对 spin 结构 $(st,V,\sigma)$。
容易看出，上述构造在相对 spin 结构的稳定共轭类与
$\widetilde{KO}(M_{[3]}/L_{[2]})$ 之间给出一一对应。
因此 $H^2(M,L;\mathbb Z_2)$ 在
$\operatorname{Spin}(M,L;\mathcal T)$ 上的作用是单传递的。

为证明 (2)，先注意：取公共细分以后，只须考虑
$\mathcal T'$ 是 $\mathcal T$ 的一个细分的情形。
此时限制给出映射
\[
\operatorname{Spin}(M,L;\mathcal T')
\longrightarrow\operatorname{Spin}(M,L;\mathcal T).
\]
依构造，该映射与 (1) 中定义的作用相容，所以它是同构。
\end{proof}

命题~\ref{prop:8-1-6} 说明，相对 spin 结构之差由
$H^2(M,L;\mathbb Z_2)$ 中的元素 $\mathfrak x$ 度量。
当用 $\mathfrak x$ 改变相对 spin 结构时，
$D_w\bar\partial$ 的指标定向改变一个因子
\[
(-1)^{\mathfrak x[w]}.
\]
（在 spin 情形，下一小节给出的例子说明了这一现象。）
以下以 $\operatorname{Spin}(M,L)$ 代替
$\operatorname{Spin}(M,L;\mathcal T)$。

由于 $st$ 由向量丛 $V$ 确定，可以略去 $st$，把相对 spin 结构
记作 $(V,\sigma)$。对辛微分同胚
\[
\psi\colon(M,L)\longrightarrow(M',L'),
\]
直接拉回给出
\[
\begin{split}
\psi^*&\colon\operatorname{Spin}(M',L')
\longrightarrow\operatorname{Spin}(M,L),\\
\psi^*[V',\sigma']&=[\psi^*V',\psi^*\sigma'].
\end{split}
\]
于是映射
\[
\psi_*\colon
\mathcal M(\beta;M,L;J)^{\psi^*[V',\sigma']}
\longrightarrow
\mathcal M(\psi_*\beta;M',L';\psi_*J)^{[V',\sigma']}
\]
保持定向。
